% Dokumentum konfiguráció
\documentclass[fleqn,12pt]{article}
% --------------------------------------------------
% Dokumentum adatok
% --------------------------------------------------

% Intézmény és kar
\newcommand\INTEZMENY{Budapesti Műszaki és Gazdaságtudományi Egyetem}
\newcommand\KAR{Gépészmérnöki Kar}
\newcommand\LOGO{bme-logo.jpg} % logo fájl neve a figures mappában

% Dokumentum adatok
\newcommand\CIM{Grundlagen der Mechatronik Hausaufgabe}
\newcommand\DATUM{2025/26/1}
\newcommand\NEV{Abos Gergő}
\newcommand\NEPTUN{C46M9M}
\newcommand\EMAIL{gergo.abos@gmail.com}
\newcommand\TANTARGY{Grundlagen der Elektronik}
\newcommand\TANTARGYKOD{BMEVIAUBSXA060-00}

\usepackage{circuitikz}
\usepackage{adjustbox}

% --------------------------------------------------
% Preambulum és egyéni parancsok betöltése
% --------------------------------------------------

\input{../_settings/00_preambulum.tex}
\input{../_settings/01_commands.tex}
\input{../_settings/02_mechTikz.tex}

% ==================================================
% Dokumentum kezdete
% ==================================================
\begin{document}

% Címlap és tartalomjegyzék
% \input{contents/01_title.tex}
\includepdf[pages=1]{contents/Aufgabenblatt.pdf} % feladatlap beszúrása, ha szükséges
\section{Bezeichnungen}
Im folgenden Dokument werde ich die Umwandlungsschritte darstellen, jeder mit seinem eigenem Schrittnummer, und die dazu gehörende Abbildung darunter. Die zu ersetzende Widerstände werden mit \textbf{Rot} bezeichnet, die neue Widerstände mit \textbf{Blau}. Die berechnungen zu jedem Schritt sind im Excel zu finden.

Erstmal ist die Aufgabe-Netz dagestellt:

\begin{figure}[H]
  \centering
  \begin{adjustbox}{width=\textwidth}
    \begin{circuitikz}[
        dot_label/.style={
          draw, circle,
          inner sep=0.5pt, minimum size=1.5em,
          fill=white,
          font=\small\bfseries
        }
      ]
      \coordinate (P1) at (0, 6);
      \coordinate (P2) at (6, 6);
      \coordinate (P3) at (12, 6);
      \coordinate (P4) at (18, 6);
      \coordinate (P5) at (0, 0);
      \coordinate (P6) at (6, 0);
      \coordinate (P7) at (12, 0);
      \coordinate (P8) at (18, 0);
      \coordinate (P9) at (20, 8);

      \coordinate (C89) at (20, 0);
      \coordinate (C39) at (12, 8);

      \coordinate (C161) at (-2, 6);
      \coordinate (C162) at (-2, 0);
      \coordinate (C163) at (-2, -2);
      \coordinate (C164) at (6, -2);

      \draw (P1) to[R, l=$R_{12}$] (P2);
      \draw (P2) to[R, l=$R_{23}$] (P3);
      \draw (P3) to[R, l=$R_{34}$] (P4);
      \draw (P5) to[R, l=$R_{56}$] (P6);
      \draw (P6) to[R, l=$R_{67}$] (P7);
      \draw (P7) to[R, l=$R_{78}$] (P8);
      \draw (P1) to[R, l=$R_{15}$] (P5);
      \draw (P2) to[R, l=$R_{26}$] (P6);
      \draw (P3) to[R, l=$R_{37}$] (P7);
      \draw (P4) to[R, l=$R_{48}$] (P8);
      \draw (P9) to[R, l=$R_{89}$] (C89);
      \draw (P9) to[R, l=$R_{39}$] (C39);
      \draw (P5) to[R, l=$R_{25}$] (P2);
      \draw (P2) to[R, l=$R_{27}$] (P7);
      \draw (P7) to[R, l=$R_{47}$] (P4);
      \draw (C161) to[R, l=$R_{16}$] (C162);
      \draw (P8) to[] (C89);
      \draw (P3) to[] (C39);
      \draw (P1) to[] (C161);
      \draw (C162) to[] (C163);
      \draw (C163) to[] (C164);
      \draw (C164) to[] (P6);

      \node[dot_label] at (P1) {1};
      \node[dot_label] at (P2) {2};
      \node[dot_label] at (P3) {3};
      \node[dot_label] at (P4) {4};
      \node[dot_label] at (P5) {5};
      \node[dot_label] at (P6) {6};
      \node[dot_label] at (P7) {7};
      \node[dot_label] at (P8) {8};
      \node[dot_label] at (P9) {9};

    \end{circuitikz}
  \end{adjustbox}
  \caption{Aufgabe-Netz}
\end{figure}


\setcounter{section}{0}
\section{Reihenschaltung vereinfachen}
\input{figures/umw_1.tex}
\section{Stern-Dreieck-Umwandlung}
\begin{figure}[H]
  \centering
  \begin{adjustbox}{width=\textwidth}
    \begin{circuitikz}[
        dot_label/.style={
          draw, circle,
          inner sep=0.5pt, minimum size=1.5em,
          fill=white,
          font=\small\bfseries
        }
      ]
      \coordinate (P1) at (0, 6);
      \coordinate (P2) at (6, 6);
      \coordinate (P3) at (12, 6);
      \coordinate (P4) at (18, 6);
      \coordinate (P5) at (0, 0);
      \coordinate (P6) at (6, 0);
      \coordinate (P7) at (12, 0);
      \coordinate (P8) at (18, 0);

      \coordinate (C89) at (20, 0);
      \coordinate (C39) at (12, 8);

      \coordinate (C161) at (-2, 6);
      \coordinate (C162) at (-2, 0);
      \coordinate (C163) at (-2, -2);
      \coordinate (C164) at (6, -2);

      \coordinate (C381) at (19, 0);
      \coordinate (C382) at (19, 2);
      \coordinate (C383) at (19, 7);
      \coordinate (C384) at (12, 7);

      \coordinate(NEW1) at (16, 4.5);
      \coordinate(NEW2) at (18, -1);
      \coordinate(NEW3) at (12, -1);
      \coordinate(NEW4) at (13, 3);

      \draw (P1) to[R, l=$R_{12}$] (P2);
      \draw (P2) to[R, l=$R_{23}$] (P3);
      \draw[red] (P3) to[R, l=$R_{34}$] (P4);
      \draw (P5) to[R, l=$R_{56}$] (P6);
      \draw (P6) to[R, l=$R_{67}$] (P7);
      \draw (P7) to[R, l=$R_{78}$] (P8);
      \draw (P1) to[R, l=$R_{15}$] (P5);
      \draw (P2) to[R, l=$R_{26}$] (P6);
      \draw (P3) to[R, l=$R_{37}$] (P7);
      \draw[red] (P4) to[R, l=$R_{48}$] (P8);
      \draw (P5) to[R, l=$R_{25}$] (P2);
      \draw (P2) to[R, l=$R_{27}$] (P7);
      \draw[red] (P7) to[R, l=$R_{47}$] (P4);
      \draw (C161) to[R, l=$R_{16}$] (C162);
      \draw (P1) to[] (C161);
      \draw (C162) to[] (C163);
      \draw (C163) to[] (C164);
      \draw (C164) to[] (P6);

      \draw (P8) to[] (C381);
      \draw (C381) to[] (C382);
      \draw (C382) to[R, l=$R_{38}$] (C383);
      \draw (C383) to[] (C384);
      \draw (C384) to[] (P3);

      \draw[blue] (P8) to[R, l=$R_{38A}$] (NEW1);
      \draw[blue] (P3) to[] (NEW1);

      \draw[blue] (P8) to[] (NEW2);
      \draw[blue] (P7) to[] (NEW3);
      \draw[blue] (NEW2) to[R, l=$R_{78A}$] (NEW3);

      \draw[blue] (P7) to[] (NEW4);
      \draw[blue] (P3) to[R, l=$R_{37A}$] (NEW4);

      \node[dot_label] at (P1) {1};
      \node[dot_label] at (P2) {2};
      \node[dot_label] at (P3) {3};
      \node[dot_label] at (P4) {4};
      \node[dot_label] at (P5) {5};
      \node[dot_label] at (P6) {6};
      \node[dot_label] at (P7) {7};
      \node[dot_label] at (P8) {8};
    \end{circuitikz}
  \end{adjustbox}
\end{figure}

\section{Parallelschaltungen vereinfachen}
\begin{figure}[H]
  \centering
  \begin{adjustbox}{width=\textwidth}
    \begin{circuitikz}[
        dot_label/.style={
          draw, circle,
          inner sep=0.5pt, minimum size=1.5em,
          fill=white,
          font=\small\bfseries
        }
      ]
      \coordinate (P1) at (0, 6);
      \coordinate (P2) at (6, 6);
      \coordinate (P3) at (12, 6);
      \coordinate (P5) at (0, 0);
      \coordinate (P6) at (6, 0);
      \coordinate (P7) at (12, 0);
      \coordinate (P8) at (18, 0);

      \coordinate (C89) at (20, 0);
      \coordinate (C39) at (12, 8);

      \coordinate (C161) at (-2, 6);
      \coordinate (C162) at (-2, 0);
      \coordinate (C163) at (-2, -2);
      \coordinate (C164) at (6, -2);

      \coordinate (C381) at (19, 0);
      \coordinate (C382) at (19, 2);
      \coordinate (C383) at (19, 7);
      \coordinate (C384) at (12, 7);
      \coordinate(C385) at (18, 6);

      \coordinate(C781) at (18, -1);
      \coordinate(C782) at (12, -1);
      \coordinate(C371) at (13, 3);

      \coordinate(NEW1) at (14, 2);
      \coordinate(NEW2) at (11, 1.5);

      \draw (P1) to[R, l=$R_{12}$] (P2);
      \draw (P2) to[R, l=$R_{23}$] (P3);
      \draw (P5) to[R, l=$R_{56}$] (P6);
      \draw (P6) to[R, l=$R_{67}$] (P7);
      \draw[red] (P7) to[R, l=$R_{78}$] (P8);
      \draw (P1) to[R, l=$R_{15}$] (P5);
      \draw (P2) to[R, l=$R_{26}$] (P6);
      \draw[red] (P3) to[R, l=$R_{37}$] (P7);
      \draw (P5) to[R, l=$R_{25}$] (P2);
      \draw (P2) to[R, l=$R_{27}$] (P7);
      \draw (C161) to[R, l=$R_{16}$] (C162);
      \draw (P1) to[] (C161);
      \draw (C162) to[] (C163);
      \draw (C163) to[] (C164);
      \draw (C164) to[] (P6);

      \draw[red] (P8) to[] (C381);
      \draw[red] (C381) to[] (C382);
      \draw[red] (C382) to[R, l=$R_{38}$] (C383);
      \draw[red] (C383) to[] (C384);
      \draw[red] (C384) to[] (P3);

      \draw[red] (P8) to[R, l=$R_{38A}$] (C385);
      \draw[red] (P3) to[] (C385);

      \draw[red] (P8) to[] (C781);
      \draw[red] (P7) to[] (C782);
      \draw[red] (C781) to[R, l=$R_{78A}$] (C782);

      \draw[red] (P7) to[] (C371);
      \draw[red] (P3) to[R, l=$R_{37A}$] (C371);

      \draw[blue] (P3) to[R, l=$R_{38B}$] (P8);

      \draw[blue] (P8) to[R, l=$R_{78B}$] (NEW1);
      \draw[blue] (P7) to[] (NEW1);

      \draw[blue] (P3) to[R, l_=$R_{37B}$] (NEW2);
      \draw[blue] (P7) to[] (NEW2);

      \node[dot_label] at (P1) {1};
      \node[dot_label] at (P2) {2};
      \node[dot_label] at (P3) {3};
      \node[dot_label] at (P5) {5};
      \node[dot_label] at (P6) {6};
      \node[dot_label] at (P7) {7};
      \node[dot_label] at (P8) {8};
    \end{circuitikz}
  \end{adjustbox}
\end{figure}

\section{Reihenschaltung vereinfachen}
\begin{figure}[H]
  \centering
  \begin{adjustbox}{width=\textwidth}
    \begin{circuitikz}[
        dot_label/.style={
          draw, circle,
          inner sep=0.5pt, minimum size=1.5em,
          fill=white,
          font=\small\bfseries
        }
      ]

      \coordinate (C38B) at (18, 6);

      \coordinate (NEW1) at (13.5, 1.5);
      \coordinate (NEW2) at (13.5, 4.5);

      \draw (P1) to[R, l=$R_{12}$] (P2);
      \draw (P2) to[R, l=$R_{23}$] (P3);
      \draw (P5) to[R, l=$R_{56}$] (P6);
      \draw (P6) to[R, l=$R_{67}$] (P7);
      \draw[red] (P7) to[R, l=$R_{78B}$] (P8);
      \draw (P1) to[R, l=$R_{15}$] (P5);
      \draw (P2) to[R, l=$R_{26}$] (P6);
      \draw (P3) to[R, l=$R_{37B}$] (P7);
      \draw (P5) to[R, l=$R_{25}$] (P2);
      \draw (P2) to[R, l=$R_{27}$] (P7);
      \draw[red] (C38B) to[] (P3);
      \draw[red] (C38B) to[R, l=$R_{38B}$] (P8);
      \draw (C161) to[R, l=$R_{16}$] (C162);
      \draw (P1) to[] (C161);
      \draw (C162) to[] (C163);
      \draw (C163) to[] (C164);
      \draw (C164) to[] (P6);

      \draw[blue] (NEW2) to[R, l=$R_{37C}$] (NEW1);
      \draw[blue] (NEW1) to[] (P7);
      \draw[blue] (NEW2) to[] (P3);

      \node[dot_label] at (P1) {1};
      \node[dot_label] at (P2) {2};
      \node[dot_label] at (P3) {3};
      \node[dot_label] at (P5) {5};
      \node[dot_label] at (P6) {6};
      \node[dot_label] at (P7) {7};
      \node[dot_label] at (P8) {8};
    \end{circuitikz}
  \end{adjustbox}
\end{figure}

\section{Parallelschaltung vereinfachen}
\begin{figure}[H]
  \centering
  \begin{adjustbox}{width=\textwidth}
    \begin{circuitikz}[
        dot_label/.style={
          draw, circle,
          inner sep=0.5pt, minimum size=1.5em,
          fill=white,
          font=\small\bfseries
        }
      ]

      \coordinate (C37C1) at (13.5, 1.5);
      \coordinate (C37C2) at (13.5, 4.5);

      \coordinate (NEW1) at (16, 0);
      \coordinate (NEW2) at (16, 6);

      \draw (P1) to[R, l=$R_{12}$] (P2);
      \draw (P2) to[R, l=$R_{23}$] (P3);
      \draw (P5) to[R, l=$R_{56}$] (P6);
      \draw (P6) to[R, l=$R_{67}$] (P7);
      \draw (P1) to[R, l=$R_{15}$] (P5);
      \draw (P2) to[R, l=$R_{26}$] (P6);
      \draw[red] (P3) to[R, l=$R_{37B}$] (P7);
      \draw (P5) to[R, l=$R_{25}$] (P2);
      \draw (P2) to[R, l=$R_{27}$] (P7);
      \draw (C161) to[R, l=$R_{16}$] (C162);
      \draw (P1) to[] (C161);
      \draw (C162) to[] (C163);
      \draw (C163) to[] (C164);
      \draw (C164) to[] (P6);

      \draw[red] (C37C2) to[R, l=$R_{37C}$] (C37C1);
      \draw[red] (C37C1) to[] (P7);
      \draw[red] (C37C2) to[] (P3);

      \draw[blue] (P7) to[] (NEW1);
      \draw[blue] (P3) to[] (NEW2);
      \draw[blue] (NEW2) to[R, l=$R_{37D}$] (NEW1);

      \node[dot_label] at (P1) {1};
      \node[dot_label] at (P2) {2};
      \node[dot_label] at (P3) {3};
      \node[dot_label] at (P5) {5};
      \node[dot_label] at (P6) {6};
      \node[dot_label] at (P7) {7};
    \end{circuitikz}
  \end{adjustbox}
\end{figure}

\section{Reihenschaltung vereinfachen}
\begin{figure}[H]
  \centering
  \begin{adjustbox}{width=\textwidth}
    \begin{circuitikz}[
        dot_label/.style={
          draw, circle,
          inner sep=0.5pt, minimum size=1.5em,
          fill=white,
          font=\small\bfseries
        }
      ]

      \coordinate (NEW1) at (11, 4);

      \draw (P1) to[R, l=$R_{12}$] (P2);
      \draw[red] (P2) to[R, l=$R_{23}$] (P3);
      \draw (P5) to[R, l=$R_{56}$] (P6);
      \draw (P6) to[R, l=$R_{67}$] (P7);
      \draw (P1) to[R, l=$R_{15}$] (P5);
      \draw (P2) to[R, l=$R_{26}$] (P6);
      \draw[red] (P3) to[R, l=$R_{37D}$] (P7);
      \draw (P5) to[R, l=$R_{25}$] (P2);
      \draw (P2) to[R, l=$R_{27}$] (P7);
      \draw (C161) to[R, l=$R_{16}$] (C162);
      \draw (P1) to[] (C161);
      \draw (C162) to[] (C163);
      \draw (C163) to[] (C164);
      \draw (C164) to[] (P6);

      \draw[blue] (P7) to[] (NEW1);
      \draw[blue] (P2) to[R, l=$R_{27A}$] (NEW1);

      \node[dot_label] at (P1) {1};
      \node[dot_label] at (P2) {2};
      \node[dot_label] at (P3) {3};
      \node[dot_label] at (P5) {5};
      \node[dot_label] at (P6) {6};
      \node[dot_label] at (P7) {7};
    \end{circuitikz}
  \end{adjustbox}
\end{figure}

\section{Parallelschaltung vereinfachen}
\begin{figure}[H]
  \centering
  \begin{adjustbox}{width=\textwidth}
    \begin{circuitikz}[
        dot_label/.style={
          draw, circle,
          inner sep=0.5pt, minimum size=1.5em,
          fill=white,
          font=\small\bfseries
        }
      ]

      \coordinate (C27A1) at (11, 4);

      \draw (P1) to[R, l=$R_{12}$] (P2);
      \draw (P5) to[R, l=$R_{56}$] (P6);
      \draw (P6) to[R, l=$R_{67}$] (P7);
      \draw (P1) to[R, l=$R_{15}$] (P5);
      \draw (P2) to[R, l=$R_{26}$] (P6);
      \draw (P5) to[R, l=$R_{25}$] (P2);
      \draw[red] (P2) to[R, l=$R_{27}$] (P7);
      \draw (C161) to[R, l=$R_{16}$] (C162);
      \draw (P1) to[] (C161);
      \draw (C162) to[] (C163);
      \draw (C163) to[] (C164);
      \draw (C164) to[] (P6);

      \draw[red] (P7) to[] (C27A1);
      \draw[red] (P2) to[R, l=$R_{27A}$] (C27A1);

      \draw[blue] (P2) to[R, l=$R_{27B}$] (P3);
      \draw[blue] (P7) to[] (P3);

      \node[dot_label] at (P1) {1};
      \node[dot_label] at (P2) {2};
      \node[dot_label] at (P5) {5};
      \node[dot_label] at (P6) {6};
      \node[dot_label] at (P7) {7};
    \end{circuitikz}
  \end{adjustbox}
\end{figure}

\section{Reihenschaltung vereinfachen}
\begin{figure}[H]
  \centering
  \begin{adjustbox}{width=\textwidth}
    \begin{circuitikz}[
        dot_label/.style={
          draw, circle,
          inner sep=0.5pt, minimum size=1.5em,
          fill=white,
          font=\small\bfseries
        }
      ]

      \coordinate (NEW1) at (7.5, 6);
      \coordinate (NEW2) at (7.5, 0);

      \draw (P1) to[R, l=$R_{12}$] (P2);
      \draw (P5) to[R, l=$R_{56}$] (P6);
      \draw[red] (P6) to[R, l=$R_{67}$] (P7);
      \draw (P1) to[R, l=$R_{15}$] (P5);
      \draw (P2) to[R, l=$R_{26}$] (P6);
      \draw (P5) to[R, l=$R_{25}$] (P2);
      \draw (C161) to[R, l=$R_{16}$] (C162);
      \draw (P1) to[] (C161);
      \draw (C162) to[] (C163);
      \draw (C163) to[] (C164);
      \draw (C164) to[] (P6);

      \draw[red] (P2) to[R, l=$R_{27B}$] (P3);
      \draw[red] (P7) to[] (P3);

      \draw[blue] (NEW1) to[R, l=$R_{26A}$] (NEW2);
      \draw[blue] (P2) to[] (NEW1);
      \draw[blue] (P6) to[] (NEW2);

      \node[dot_label] at (P1) {1};
      \node[dot_label] at (P2) {2};
      \node[dot_label] at (P5) {5};
      \node[dot_label] at (P6) {6};
      \node[dot_label] at (P7) {7};
    \end{circuitikz}
  \end{adjustbox}
\end{figure}

\section{Parallelschaltung vereinfachen}
\input{figures/umw_9.tex}
\section{Stern-Dreieck-Umwandlung}
\begin{figure}[H]
  \centering
  \begin{adjustbox}{width=\textwidth}
    \begin{circuitikz}[
        dot_label/.style={
          draw, circle,
          inner sep=0.5pt, minimum size=1.5em,
          fill=white,
          font=\small\bfseries
        }
      ]

      \coordinate (NEW1) at (7.5, 0);
      \coordinate (NEW2) at (7.5, 6);
      \coordinate (NEW3) at (0, 7.5);
      \coordinate (NEW4) at (6, 7.5);
      \coordinate (NEW5) at (5, 4);

      \draw (P1) to[R, l=$R_{12}$] (P2);
      \draw[red] (P5) to[R, l=$R_{56}$] (P6);
      \draw[red] (P1) to[R, l=$R_{15}$] (P5);
      \draw (P2) to[R, l=$R_{26B}$] (P6);
      \draw[red] (P5) to[R, l=$R_{25}$] (P2);
      \draw (C161) to[R, l=$R_{16}$] (C162);
      \draw (P1) to[] (C161);
      \draw (C162) to[] (C163);
      \draw (C163) to[] (C164);
      \draw (C164) to[] (P6);

      \draw[blue] (NEW2) to[R, l=$R_{26C}$] (NEW1);
      \draw[blue] (P6) to[] (NEW1);
      \draw[blue] (NEW2) to[] (P2);

      \draw[blue] (NEW3) to[R, l=$R_{12A}$] (NEW4);
      \draw[blue] (P1) to[] (NEW3);
      \draw[blue] (NEW4) to[] (P2);

      \draw[blue] (P6) to[] (NEW5);
      \draw[blue] (P1) to[R, l=$R_{16A}$] (NEW5);

      \node[dot_label] at (P1) {1};
      \node[dot_label] at (P2) {2};
      \node[dot_label] at (P5) {5};
      \node[dot_label] at (P6) {6};
    \end{circuitikz}
  \end{adjustbox}
\end{figure}

\section{Parallelschaltungen vereinfachen}
\begin{figure}[H]
  \centering
  \begin{adjustbox}{width=\textwidth}
    \begin{circuitikz}[
        dot_label/.style={
          draw, circle,
          inner sep=0.5pt, minimum size=1.5em,
          fill=white,
          font=\small\bfseries
        }
      ]

      \coordinate (C26C1) at (7.5, 0);
      \coordinate (C26C2) at (7.5, 6);
      \coordinate (C12A1) at (0, 7.5);
      \coordinate (C12A2) at (6, 7.5);
      \coordinate (C16A1) at (0, 0);

      \coordinate (NEW1) at (9, 0);
      \coordinate (NEW2) at (9, 6);

      \draw[blue] (P1) to[R, l=$R_{16B}$] (P6);
      \draw[blue] (NEW2) to[R, l=$R_{26D}$] (NEW1);
      \draw[blue] (NEW2) to[] (P2);
      \draw[blue] (P6) to[] (NEW1);

      \draw (P1) to[R, l=$R_{12}$] (P2);
      \draw[red] (P2) to[R, l=$R_{26B}$] (P6);
      \draw[red] (C161) to[R, l=$R_{16}$] (C162);
      \draw[red] (P1) to[] (C161);
      \draw[red] (C162) to[] (C163);
      \draw[red] (C163) to[] (C164);
      \draw[red] (C164) to[] (P6);

      \draw[red] (C26C2) to[R, l=$R_{26C}$] (C26C1);
      \draw[red] (P6) to[] (C26C1);
      \draw[red] (C26C2) to[] (P2);

      \draw (C12A1) to[R, l=$R_{12A}$] (C12A2);
      \draw (P1) to[] (C12A1);
      \draw (C12A2) to[] (P2);

      \draw[red] (P6) to[] (C16A1);
      \draw[red] (P1) to[R, l=$R_{16A}$] (C16A1);

      \node[dot_label] at (P1) {1};
      \node[dot_label] at (P2) {2};
      \node[dot_label] at (P6) {6};
    \end{circuitikz}
  \end{adjustbox}
\end{figure}

\section{Reihenschaltung vereinfachen}
\begin{figure}[H]
  \centering
  \begin{adjustbox}{width=\textwidth}
    \begin{circuitikz}[
        dot_label/.style={
          draw, circle,
          inner sep=0.5pt, minimum size=1.5em,
          fill=white,
          font=\small\bfseries
        }
      ]

      \coordinate (NEW1) at (6, 5);
      \coordinate (NEW2) at (2, 5);

      \draw[red] (P1) to[R, l=$R_{16B}$] (P6);
      \draw[red] (P6) to[R, l=$R_{26D}$] (P2);

      \draw (P1) to[R, l=$R_{12}$] (P2);

      \draw (C12A1) to[R, l=$R_{12A}$] (C12A2);
      \draw (P1) to[] (C12A1);
      \draw (C12A2) to[] (P2);

      \draw[blue] (P1) to[] (NEW2);
      \draw[blue] (P2) to[] (NEW1);
      \draw[blue] (NEW2) to[R, l=$R_{12B}$] (NEW1);

      \node[dot_label] at (P1) {1};
      \node[dot_label] at (P2) {2};
      \node[dot_label] at (P6) {6};
    \end{circuitikz}
  \end{adjustbox}
\end{figure}

\section{Parallelschaltung vereinfachen}
\begin{figure}[H]
  \centering
  \begin{adjustbox}{width=\textwidth}
    \begin{circuitikz}[
        dot_label/.style={
          draw, circle,
          inner sep=0.5pt, minimum size=1.5em,
          fill=white,
          font=\small\bfseries
        }
      ]

      \coordinate (NEW1)at (0, 3);
      \coordinate (NEW2)at (6, 3);

      \draw[blue] (P1) to[] (NEW1);
      \draw[blue] (P2) to[] (NEW2);
      \draw[blue] (NEW1) to[R, l=$R_{Gesamt}$] (NEW2);

      \coordinate (C12B2) at (6, 4.5);
      \coordinate (C12B1) at (0, 4.5);

      \draw[red] (P1) to[R, l=$R_{12}$] (P2);

      \draw[red] (C12A1) to[R, l=$R_{12A}$] (C12A2);
      \draw[red] (P1) to[] (C12A1);
      \draw[red] (C12A2) to[] (P2);

      \draw[red] (P1) to[] (C12B1);
      \draw[red] (P2) to[] (C12B2);
      \draw[red] (C12B1) to[R, l=$R_{12B}$] (C12B2);

      \node[dot_label] at (P1) {1};
      \node[dot_label] at (P2) {2};
    \end{circuitikz}
  \end{adjustbox}
\end{figure}

\end{document}
